\documentclass[11pt,a4paper]{article}

\usepackage[utf8]{inputenc}
\usepackage[T1]{fontenc}
\usepackage[margin=2.5cm]{geometry}
\usepackage{graphicx}
\usepackage{booktabs}
\usepackage{xcolor}
\usepackage{hyperref}
\usepackage{listings}
\usepackage{fancyhdr}
\usepackage{tikz}
\usepackage{longtable}
\usepackage{caption}
\usepackage{enumitem}

\definecolor{codebg}{RGB}{245,245,245}
\definecolor{accent}{RGB}{28,110,140}

\lstset{
  backgroundcolor=\color{codebg},
  basicstyle=\ttfamily\scriptsize,
  breaklines=true,
  frame=single,
  rulecolor=\color{gray!40},
  columns=fullflexible,
  captionpos=b
}

\hypersetup{
  colorlinks=true,
  linkcolor=accent,
  urlcolor=accent,
  citecolor=accent
}

\pagestyle{fancy}
\fancyhf{}
\setlength{\headheight}{14pt}
\rhead{Trabajo Unidad II — BigData 2026A}
\lhead{UNSA-EPCC}
\rfoot{\thepage}

\newcommand{\placeholderCaptura}[2]{%
\begin{center}
\begin{tikzpicture}
\node[draw, dashed, gray, minimum width=13cm, minimum height=#2,
      text width=12cm, align=center, font=\ttfamily\small, color=gray]
      {[Insertar captura aqu\'i: \\ \texttt{#1}]};
\end{tikzpicture}
\end{center}
}

\title{
\vspace{-1cm}
\begin{center}
\textbf{UNIVERSIDAD NACIONAL DE SAN AGUST\'IN DE AREQUIPA}\\
\textbf{ESCUELA PROFESIONAL DE CIENCIA DE LA COMPUTACI\'ON}\\[0.2cm]
\textsc{Curso: BigData 2026A}\\[1cm]
\rule{\textwidth}{0.5pt}\\[0.3cm]
{\LARGE \textbf{Trabajo Unidad III}}\\[0.3cm]
{\Large Plataforma Inteligente para Simulaci\'on y An\'alisis de\\
Audiencias Digitales en Tiempo Real utilizando\\
Arquitecturas Orientadas a Eventos}\\[0.3cm]
\rule{\textwidth}{0.5pt}
\end{center}
}

\author{Rivas Huanca Diego Raul \\
Sumare Uscca Josue Gabriel \\
Humpiri Valdivia Freddy Leonel\\
Escuela Profesional de Ciencia de la Computaci\'on}
\date{Julio 2026}

\begin{document}

\begin{titlepage}
\maketitle
\thispagestyle{empty}
\vfill
\begin{center}
\small Repositorio del proyecto: \url{https://github.com/DiegoRivas1/bd-audiencias-digitales}
\end{center}
\end{titlepage}

\renewcommand{\contentsname}{Índice}

\tableofcontents
\newpage

% ============================================================
\section{Introducci\'on}
% ============================================================

Las organizaciones modernas generan volúmenes masivos de eventos digitales, inicios de sesión, búsquedas, visualizaciones de producto, compras que
tradicionalmente se analizaban mediante procesamiento por lotes, con
retrasos de horas o días. Este trabajo implementa una plataforma de
\textbf{procesamiento de datos en tiempo real} basada en una
\textbf{Arquitectura Orientada a Eventos (EDA)}, capaz de simular el
comportamiento de decenas de usuarios digitales mediante agentes
autónomos, capturar sus eventos con Apache Kafka, procesarlos
continuamente con Apache Flink para construir audiencias digitales, y
visualizar los resultados en un dashboard interactivo en tiempo real.

El objetivo del proyecto es demostrar el flujo completo
\emph{agentes~$\rightarrow$~Kafka~$\rightarrow$~Flink~$\rightarrow$~audiencias~$\rightarrow$~dashboard},
y analizar cómo cambia el comportamiento agregado de la plataforma bajo
distintos escenarios comerciales (día normal, Cyber Monday, campaña
escolar, etc.).

% ============================================================
\section{Marco conceptual}
% ============================================================

\subsection{Arquitectura Orientada a Eventos (EDA)}
Un modelo de diseño donde los componentes de software no se comunican
directamente entre sí, sino que publican \textbf{eventos} acciones
significativas ocurridas en el sistema en un canal común, permitiendo que
otros componentes los consuman de forma independiente y desacoplada.

\subsection{Event Streaming y Apache Kafka}
El \emph{Event Streaming} es el procesamiento continuo de eventos
conforme se producen, a diferencia del procesamiento por lotes. Apache
Kafka es la plataforma de streaming utilizada como intermediario:
organiza los eventos en \textbf{tópicos}, particionados para permitir
paralelismo y alta disponibilidad.

\subsection{Apache Flink}
Motor de procesamiento de flujos de datos (\emph{Stream Processing
Engine}) usado para consumir los eventos desde Kafka y aplicar
transformaciones continuas: filtrado, enriquecimiento, ventanas de
tiempo, agregaciones y clasificación basada en reglas de negocio.

\subsection{Audiencias digitales}
Conjuntos de usuarios agrupados según patrones de comportamiento
observados en tiempo real (y no solo datos demográficos estáticos), útiles
para personalizar campañas, detectar riesgo de abandono o identificar
clientes de alto valor.

% ============================================================
\section{Arquitectura del sistema}
% ============================================================

El sistema implementado sigue el flujo:

\begin{center}
\begin{verbatim}
Agentes simulados (Python, 8 perfiles)
        |  publican eventos JSON
        v
Apache Kafka  (audiencias-user-events, audiencias-purchase-events)
        |  consumido continuamente
        v
Apache Flink  (metricas por ventana + clasificador de audiencias)
        |  publica resultados JSON
        v
Kafka  (audiencias-metricas-output, audiencias-resultado-output)
        |  consumido por
        v
Backend FastAPI (WebSocket)  --->  Dashboard web en tiempo real
\end{verbatim}
\end{center}

A diferencia de una arquitectura con base de datos intermedia, aquí Flink
publica sus resultados directamente a nuevos tópicos de Kafka, y un
backend ligero en FastAPI los retransmite por WebSocket al navegador. Esto
evita la complejidad de una capa de persistencia adicional, manteniendo
el principio de streaming de punta a punta: cada resultado se empuja al
cliente apenas se calcula, sin necesidad de que el dashboard haga polling
a una base de datos.

\subsection{Infraestructura desplegada}
Todo el sistema corre sobre una instancia EC2 (Ubuntu 24.04, \texttt{t3.large},
2 vCPU / 8 GiB RAM) en AWS Academy, con cuatro procesos concurrentes:
el broker de Kafka, el job de Flink, el backend del dashboard, y el
simulador de agentes.

% ============================================================
\section{Dise\~no e implementaci\'on}
% ============================================================

\subsection{Simulador de agentes}

Se implementaron 8 perfiles de comportamiento, cada uno corriendo en su
propio hilo y generando sesiones completas de navegación (login
$\rightarrow$ ver/buscar productos $\rightarrow$ carrito
$\rightarrow$ compra o abandono $\rightarrow$ logout):

\begin{center}
\begin{longtable}{@{}p{3.2cm}p{9.5cm}@{}}
\toprule
\textbf{Perfil} & \textbf{Comportamiento} \\
\midrule
\endhead
Comprador compulsivo & Compra rápido, consulta pocos productos, alta probabilidad de compra \\
Comparador & Consulta muchos productos, compara precios, compra ocasionalmente \\
Comprador nocturno & Solo genera actividad entre las 22:00 y las 05:00 \\
Cliente Premium & Compra pocas veces, pero de alto valor (S/1500 -- S/6000) \\
Cliente Frecuente & Realiza compras constantemente, alta frecuencia de sesión \\
Usuario Explorador & Navega mucho (8--20 productos), nunca compra \\
Cliente Indeciso & Agrega y quita productos del carrito repetidamente \\
Cliente Estacional & Su comportamiento se amplifica según el escenario comercial activo \\
\bottomrule
\end{longtable}
\end{center}

Cada agente publica eventos JSON con la estructura:
\begin{lstlisting}
{
  "timestamp": "2026-07-20T02:21:53",
  "user_id": "USR55301",
  "company": "Retail",
  "event": "VIEW_PRODUCT",
  "product": "Laptop Lenovo",
  "category": "Electronics",
  "city": "Arequipa",
  "price": 3341.50,
  "agent_id": "CLIENTE_ESTACIONAL_01",
  "agent_type": "CLIENTE_ESTACIONAL"
}
\end{lstlisting}

Adicionalmente, se implementaron \textbf{7 escenarios comerciales}
(\texttt{normal}, \texttt{navidad}, \texttt{cyber\_monday},
\texttt{black\_friday}, \texttt{dia\_del\_padre},
\texttt{fiestas\_patrias}, \texttt{campana\_escolar}), cada uno aplicando
multiplicadores globales sobre la probabilidad de compra, la frecuencia
de sesiones y el valor promedio de compra de \emph{todos} los agentes,
simulando el efecto de una campaña real sobre el tráfico de la
plataforma.

\subsection{Apache Kafka}

Se configuraron 4 tópicos:

\begin{center}
\begin{tabular}{@{}llc@{}}
\toprule
\textbf{Tópico} & \textbf{Contenido} & \textbf{Particiones} \\
\midrule
\texttt{audiencias-user-events} & LOGIN, SEARCH, VIEW\_PRODUCT, ADD\_CART, REMOVE\_CART, LOGOUT & 6 \\
\texttt{audiencias-purchase-events} & PURCHASE, PAYMENT\_FAILED & 6 \\
\texttt{audiencias-metricas-output} & Métricas agregadas por ventana (salida de Flink) & 1 \\
\texttt{audiencias-resultado-output} & Audiencias detectadas por usuario (salida de Flink) & 3 \\
\bottomrule
\end{tabular}
\end{center}

\subsection{Apache Flink: métricas y clasificación de audiencias}

Se implementó un único job (\texttt{MainDashboardPipeline}) con dos
ramas de procesamiento sobre el mismo stream de eventos:

\textbf{Rama de métricas} (ventanas de 15 segundos, \texttt{TumblingProcessingTimeWindows}):
eventos totales, eventos por segundo, usuarios activos distintos, eventos
por tipo, top 5 productos más vistos/comprados, compras por ciudad, y
tasa de conversión (\texttt{PURCHASE} / \texttt{ADD\_CART}).

\textbf{Rama de audiencias} (estado por usuario mediante
\texttt{KeyedProcessFunction}, \texttt{keyBy(user\_id)}): mantiene
contadores de sesión (vistas, vistas en categoría Electronics, productos
agregados al carrito, carrito activo, valor máximo de compra) y un
contador histórico de compras que persiste entre sesiones. Al detectar el
evento \texttt{LOGOUT}, evalúa las siguientes reglas:

\begin{center}
\begin{longtable}{@{}p{3.6cm}p{9cm}@{}}
\toprule
\textbf{Audiencia} & \textbf{Regla} \\
\midrule
\endhead
Alta intención de compra / Abandono de carrito & Agregó productos al carrito pero no completó ninguna compra en la sesión \\
Cliente Frecuente & 3 o más compras completadas (histórico acumulado) \\
Cliente Premium & Alguna compra de la sesión superó S/1500 \\
Interesado en Tecnología & Al menos 60\% de sus vistas fueron de la categoría \emph{Electronics} \\
Riesgo de Abandono & Navegó bastante (5+ vistas) pero nunca agregó nada al carrito \\
\bottomrule
\end{longtable}
\end{center}

Un mismo usuario puede activar varias audiencias simultáneamente (por
ejemplo, \texttt{ALTA\_INTENCION\_COMPRA} \texttt{|} \texttt{INTERESADO\_TECNOLOGIA}),
lo cual se observó consistentemente en las pruebas.

\subsection{Dashboard en tiempo real}

Backend en FastAPI que consume los tópicos de salida de Flink
(\texttt{KafkaConsumer} en un hilo aparte) y retransmite cada mensaje
nuevo a todos los clientes conectados vía WebSocket, sin persistencia
intermedia. El frontend (HTML/CSS/JS + Chart.js) muestra:

\begin{itemize}[nosep]
\item KPIs en vivo: eventos totales, eventos/segundo, usuarios activos, conversión
\item Gráfico de eventos por tipo y de compras por ciudad
\item Tablas de productos más vistos y más comprados
\item Feed en vivo de audiencias detectadas por usuario
\item Alertas automáticas (ej. conversión baja, alta tasa de pagos rechazados)
\end{itemize}

% ============================================================
\section{Resultados: comparativa de escenarios}
% ============================================================

Se ejecutó el simulador (20 agentes, 8 perfiles) bajo tres escenarios de
forma secuencial y aislada (un único proceso productor activo por
corrida, confirmado mediante el conteo de usuarios activos en el
dashboard, siempre $\leq$ 20), dejando correr cada uno un par de minutos
antes de capturar resultados.

\begin{center}
\begin{tabular}{@{}lccccc@{}}
\toprule
\textbf{Escenario} & \textbf{Eventos} & \textbf{Ev./seg} & \textbf{Usuarios} & \textbf{Conversión} & \textbf{Mult. prob. compra} \\
\midrule
Campaña Escolar & 128 & 8.53 & 18 & 23.1\% & 1.3$\times$ \\
Normal (base) & 111 & 7.40 & 18 & 31.6\% & 1.0$\times$ \\
Cyber Monday & 150 & 10.00 & 18 & 33.3\% & 2.2$\times$ \\
\bottomrule
\end{tabular}
\end{center}

\subsection{Análisis}

Los resultados muestran una tendencia \textbf{coherente con la
configuración de multiplicadores} definida para cada escenario: Campaña
Escolar (multiplicador de compra más bajo, 1.3$\times$) produjo la menor
conversión observada (23.1\%), mientras que Cyber Monday (multiplicador
más alto, 2.2$\times$) produjo la mayor (33.3\%), con el escenario Normal
como punto intermedio. El volumen de eventos por segundo también creció
en el mismo orden (Campaña Escolar $<$ Normal $<$ Cyber Monday), reflejando
el multiplicador de frecuencia de sesiones de cada escenario.

Es importante notar que la conversión no escaló de forma estrictamente
lineal con el multiplicador (2.2$\times$ el multiplicador no produjo
2.2$\times$ la conversión base), lo cual es esperado: la probabilidad de
compra está acotada a un máximo de 0.95 en el motor del agente
(\texttt{perfiles.py}), y no todos los perfiles reaccionan igual al
escenario (por ejemplo, el Usuario Explorador nunca compra
independientemente del escenario activo).

En una corrida adicional de mayor duración bajo \texttt{cyber\_monday}
(8632 eventos acumulados, 20 usuarios activos), la conversión general se
estabilizó en 36.9\%, con \emph{Mouse Logitech}, \emph{Mochila Escolar} y
\emph{Laptop Lenovo} como los productos más comprados, y una distribución
de compras razonablemente uniforme entre las 5 ciudades simuladas
(Piura, Lima, Arequipa, Cusco, Trujillo), lo que confirma que el
simulador no introduce sesgos geográficos artificiales.

\subsection{Audiencias detectadas por escenario}

En los tres escenarios se observaron consistentemente las 5 audiencias
definidas, con las combinaciones \texttt{ALTA\_INTENCION\_COMPRA /
ABANDONO\_CARRITO} e \texttt{INTERESADO\_TECNOLOGIA} como las más
frecuentes esperable dado que el catálogo de productos tiene una alta
proporción de artículos de la categoría \emph{Electronics}. Se detectaron
además usuarios con múltiples audiencias simultáneas (por ejemplo,
\texttt{CLIENTE\_PREMIUM | INTERESADO\_TECNOLOGIA}), demostrando que el
clasificador no trata las audiencias como mutuamente excluyentes,
reflejando mejor el comportamiento real de un usuario.

% ============================================================
\section{Capturas de pantalla}
% ============================================================

\subsection{Dashboard — Escenario Normal}
\begin{center}
\includegraphics[width=0.9\textwidth]{docs/capturas/dashboard_normal.png}
\end{center}

\subsection{Dashboard — Escenario Campaña Escolar}
\begin{center}
\includegraphics[width=0.9\textwidth]{docs/capturas/dashboard_campana_escolar.png}
\end{center}

\subsection{Dashboard — Escenario Cyber Monday}
\begin{center}
\includegraphics[width=0.9\textwidth]{docs/capturas/dashboard_cyber_monday.png}
\end{center}



% ============================================================
\section{Conclusiones}
% ============================================================

\begin{enumerate}
\item Se implementó exitosamente el flujo completo de una arquitectura
orientada a eventos, desde la generación simulada de tráfico hasta la
visualización en tiempo real, sin componentes de procesamiento por
lotes en ninguna etapa.

\item El diseño de agentes con perfiles de comportamiento diferenciados,
combinado con escenarios que aplican multiplicadores globales, permitió
observar cambios de comportamiento agregado consistentes y explicables
por el diseño (mayor multiplicador de probabilidad de compra
$\rightarrow$ mayor conversión observada), validando que el simulador
responde correctamente a los parámetros configurados.

\item El uso de \texttt{KeyedProcessFunction} en Flink permitió mantener
estado por usuario a lo largo de una sesión (y parcialmente entre
sesiones, para el conteo histórico de compras) sin necesidad de una base
de datos externa, resolviendo la clasificación de audiencias
completamente en memoria dentro del propio job de streaming.

\item Publicar los resultados de Flink directamente a tópicos de Kafka
de salida, consumidos por un backend WebSocket sin capa de persistencia,
demostró ser una arquitectura suficiente y más simple que introducir una
base de datos intermedia, para el caso de uso de visualización en tiempo
real de este proyecto.

\item Entre las dificultades técnicas resueltas destacan: conflictos de
classloading de Flink al ejecutarse vía plugins de Maven (resuelto
empaquetando un \emph{fat jar} con \texttt{maven-shade-plugin}),
incompatibilidad de Kryo con el sistema de módulos de Java 17 al
serializar colecciones en POJOs (resuelto evitando \texttt{List} en
favor de \texttt{String} en los objetos de resultado), y conflictos de
versión entre \texttt{jackson-core} y \texttt{jackson-databind}
resueltos fijando versiones explícitas en el \texttt{pom.xml}.
\end{enumerate}

% ============================================================
\section{Referencias}
% ============================================================

\begin{itemize}[nosep]
\item Apache Software Foundation. (2026). \emph{Apache Kafka Documentation}. \url{https://kafka.apache.org/documentation/}
\item Apache Software Foundation. (2026). \emph{Apache Flink Documentation}. \url{https://nightlies.apache.org/flink/flink-docs-stable/}
\item Kleppmann, M. (2017). \emph{Designing Data-Intensive Applications}. O'Reilly Media.
\item AWS Documentation. (2026). \emph{Amazon Managed Service for Apache Flink}. \url{https://docs.aws.amazon.com/managed-flink/}
\end{itemize}

\end{document}